%!TEX root = ../../thesis.tex

\section{States and Observations}\label{sec:mdp-states}

For a given PUT, we model the state space as the joint state of process memory,
kernel state plus register file and other execution state necessary to PUT execution.
This constitutes a full description of the PUT and therefore suffices to fully predict future
behavior of the PUT given new inputs.
Although this state space formulation is straightforward, it is too large to serve as a
base for a fuzzing implementation.
In practice, this state is contained in the execution environment and even though it
is accessible to us, we will not use it directly.
Adhering to the approach of \gls{cgf}, the fuzzer instead observes coverage information $\mathcal{C}$.
In our case, this can either be block or edge coverage.
We therefore model the observation function $o_F:\mathcal{S}\rightarrow\mathcal{C}$ to map
true PUT state $\mathcal{S}_F$ to coverage feedback $\mathcal{C}$.
Depending on the action model, we extend the observation function to
include additional information.

Observations must be encoded so that the representation network $h$ can properly
consume them.
The block hit tracker has a predetermined number of buckets, one for
each coverage input dimension of $h$.
For a given bucket, we convert its value from an integer to floating point and
feed this to the representation network, allowing it to encode which
blocks or edges have been hit and how often.
Another valid approach for smaller binaries lies in bucketing the hit count and then
$k$-hot encoding these buckets.
This is in theory a more correct approach, but causes an increase in the input dimensions of $h$
proportional to the number of hit count buckets.
